% Part: lambda-calculus
% Chapter: syntax
% Section: de-bruijn

\documentclass[../../../include/open-logic-section]{subfiles}

\begin{document}

\olfileid{lam}{syn}{deb}

\olsection{డి బ్రూయిన్ సూచిక}

$\alpha$-తుల్యత సహజమైన భావన: $\alpha$-తుల్య పదాలకు ఒకే
అర్థం ఉంటుంది. ఒకదానిగా మరొకటి $\alpha$-పరివర్తనం చెందగల
పదాలను వేర్వేరుగా చూడని లాంబ్డా పదాల వాక్యనిర్మాణాన్ని కూడా
ఇవ్వవచ్చు. అందులో ప్రసిద్ధమైన పద్ధతి చరాల స్థానంలో వాటి
\emph{డి బ్రూయిన్ సూచికలను} వాడుతుంది.

$\lambd[x][M]$ అని రాస్తే, ప్రమేయపు పరామితి~$x$ అని స్పష్టంగా
చెబుతున్నాం; కాబట్టి $M$లో $x$తో ఆ పరామితిని సూచించవచ్చు.
డి బ్రూయిన్ సూచికలలో మాత్రం బద్ధ పరామితులకు పేర్లు ఉండవు.
ప్రమేయపు అంతర్భాగంలో వాటిని సూచించడానికి, చరం ఘటనకూ దానిని
బంధించే అమూర్తీకరణకూ మధ్యలో ఉన్న అమూర్తీకరణల సంఖ్యను వాడుతాం.
ఉదాహరణకు $\lambd[x][\lambd[y][y x]]$ను చూడండి: బాహ్య
అమూర్తీకరణ~$x$ను, అంతర అమూర్తీకరణ~$y$ను బంధిస్తుంది.
$y x$ ఉపపదం అంతర అమూర్తీకరణ పరిధిలో ఉంది. $y$ ఘటనకూ
దాని బంధకం~$\lambd[y]$కూ మధ్యలో అమూర్తీకరణ లేదు; $x$ ఘటనకూ
దాని బంధకం~$\lambd[x]$కూ మధ్యలో ఒకటి ఉంది. కనుక
$y x$ను $0\, 1$గా, మొత్తం పదాన్ని
$\lambd[][\lambd[][0\,1]]$గా రాస్తాం.
\sourcecorrection{OLTELAMDEB-001}{మూలంలోని ఉదాహరణ వివరణలో “outer abstraction is on binds”, “binds the variable is” వంటి ఆంగ్ల వ్యాకరణ దోషాలున్నాయి; తెలుగు వాక్యంలో బాహ్య అమూర్తీకరణ xను, అంతర అమూర్తీకరణ yను బంధిస్తాయని స్పష్టంగా ఇచ్చాం. మూలం ఉపపదాన్ని $0\,1$గా, మొత్తం పదంలో అదే ఉపపదాన్ని $01$గా రాస్తుంది; రెండు సూచికల ప్రయోగం స్పష్టమయ్యేలా మొత్తం పదంలోనూ $0\,1$గా చూపాం.}

\begin{defn}
  డి బ్రూయిన్ పదాలను ఆగమన పద్ధతిలో ఇలా నిర్వచిస్తాం:
  \begin{enumerate}
  \item $n$, ఇక్కడ $n$ ఏ సహజ సంఖ్యైనా కావచ్చు.
  \item $PQ$, ఇక్కడ $P$, $Q$ రెండూ డి బ్రూయిన్ పదాలు.
  \item $\lambd[][N]$, ఇక్కడ $N$ డి బ్రూయిన్ పదం.
  \end{enumerate}
\end{defn}

సాధారణ లాంబ్డా పదాల నుంచి డి బ్రూయిన్ సూచికలతో కూడిన పదాలకు
ఒక ఆచారబద్ధ అనువాదం ఇలా ఉంటుంది:
\begin{defn}
  \begin{align*}
    F_\Gamma(x) &= \Gamma(x) \\
    F_\Gamma(PQ) &= F_\Gamma(P)F_\Gamma(Q) \\
    F_\Gamma(\lambd[x][N]) &= \lambd[][F_{x,\Gamma}(N)]
  \end{align*}
  ఇక్కడ $\Gamma$ సున్నా నుంచి సూచికలిచ్చిన చరాల జాబితా;
  $\Gamma(x)$ అనేది $\Gamma$ జాబితాలో చరం~$x$ యొక్క స్థానం.
  ఉదాహరణకు $\Gamma$ అనేది $x,y,z$ అయితే,
  $\Gamma(x)$ అనేది $0$; $\Gamma(z)$ అనేది $2$.
  ఒకే చరం జాబితాలో పలుసార్లు ఉంటే, బంధనానికి దగ్గరైన
  మొదటి ఘటన స్థానం తీసుకుంటాం.
  \sourcecorrection{OLTELAMDEB-002}{మూలం $\Gamma(x)$ను “జాబితాలో $x$ స్థానం”గా మాత్రమే వివరిస్తుంది. $x,\Gamma$ నిర్మాణం ఒకే చరాన్ని మళ్లీ జాబితా తలలో ఉంచవచ్చు; అప్పుడు స్థానాన్ని నిర్దిష్టం చేయడానికి మొదటి (కనిష్ఠ సూచిక) ఘటనను తీసుకోవాలి. ఇది లోపలి బంధకం బాహ్యదాన్ని కప్పే సాధారణ బంధన నియమంతో సరిపోతుంది.}

  $x,\Gamma$ అంటే $\Gamma$ తలలో $x$ను చేర్చి పొందిన జాబితా.
  ఉదాహరణకు పైన ఉన్న $\Gamma$కే కొనసాగింపుగా
  $w,\Gamma$ అనేది $w,x,y,z$.
\end{defn}

డి బ్రూయిన్ పదం నుంచి సాధారణ లాంబ్డా పదాన్ని తిరిగి ఇలా పొందుతాం:

\begin{defn}
  \begin{align*}
    G_\Gamma(n) &= \Gamma[n] \\
    G_\Gamma(PQ) &= G_\Gamma(P) G_\Gamma(Q) \\
    G_\Gamma(\lambd[][N]) &= \lambd[x][G_{x,\Gamma}(N)]
  \end{align*}
  ఇక్కడ కూడా $\Gamma$ సున్నా నుంచి సూచికలిచ్చిన చరాల జాబితా;
  $\Gamma[n]$ అనేది $n$వ స్థానంలోని చరం. ఉదాహరణకు
  $\Gamma$ అనేది $x,y,z$ అయితే $\Gamma[1]$ అనేది $y$.

  చివరి సమీకరణంలో $x$ను $\Gamma$లో లేని ఏ చరంగానైనా
  ఎంచుకుంటాం. అలాగే $n$ జాబితా పరిధిలో లేని చోట
  $G_\Gamma(n)$ నిర్వచితం కాదు; ఈ తిరుగు అనువాదం
  సంబంధిత జాబితాకు సరిగ్గా సూచికలున్న పదాలకే వర్తిస్తుంది.
  \sourcecorrection{OLTELAMDEB-003}{మూలం ప్రతి సహజ సంఖ్యను డి బ్రూయిన్ పదంగా అనుమతిస్తుంది, కానీ $G_\Gamma(n)=\Gamma[n]$ను అన్ని $n$లకు నిర్వచించదు. జాబితా పొడవు మించిన సూచికలకు తిరుగు పటం పాక్షికమని వెల్లడించాం; మూల సమీకరణాన్ని మార్చలేదు.}
\end{defn}

నిరూపణలు ఇవ్వకుండా కొన్ని ఫలితాలను ఇక్కడ పేర్కొంటాం:

\begin{prop}
  $M \aconvone M'$, $\Gamma$ అనేది $\FV{M}$లోని
  చరాలన్నీ ఉన్న ఏ జాబితా అయినా, అప్పుడు
  $F_\Gamma(M) \eqs F_\Gamma(M')$.
\end{prop}

\end{document}
